Este trabalho teve como objetivo entender os desafios de atualizar sistemas legados com agentes de IA, do ponto de vista do desenvolvimento.

Na literatura, encontramos desafios em seis pontos: tamanho do sistema, complexidade das funcionalidades, regras escondidas, atalhos do agente, linguagens antigas e custo. Em todos os estudos, o resultado do agente precisou ser verificado, e os testes foram a forma mais usada para isso.

A proposta do estudo de caso mostra como observar esses desafios em um sistema real: usar o agente em partes pequenas, proteger os testes e registrar cada problema encontrado.

% TODO: depois do estudo de caso, escrever aqui o que foi encontrado:
% quais desafios apareceram, quais não apareceram e o que o desenvolvedor pode fazer.

Como trabalhos futuros, o mesmo método pode ser aplicado em outros sistemas, com outros agentes e outras trocas de tecnologia, para comparar os resultados.
